%======================================================================
% چکیده
%======================================================================
\chapter*{چکیده}
\addcontentsline{toc}{chapter}{چکیده}

با گسترش روزافزون \term{سرویس‌های هوش مصنوعی}{AI Services} و تنوع
\term{رابط‌های برنامه‌نویسی کاربردی}{Application Programming Interface (API)}
ارائه‌شده توسط شرکت‌های مختلف، سازمان‌ها و توسعه‌دهندگان با چالش
یکپارچه‌سازی، مدیریت دسترسی و صورت‌حساب‌گذاری این سرویس‌ها روبه‌رو هستند.
این پایان‌نامه طراحی و پیاده‌سازی یک \term{سامانه‌ی واسط}{Gateway} یکپارچه
برای دسترسی به سرویس‌های هوش مصنوعی را شرح می‌دهد.

سامانه‌ی پیاده‌سازی‌شده با استفاده از چارچوب \lr{Laravel} (زبان \lr{PHP}) در
سمت سرور، پایگاه داده‌ی \lr{SQL Server}، سیستم \term{کش}{Cache} و
\term{صف}{Queue} مبتنی بر \lr{Redis}، و رابط کاربری \lr{Vue.js} طراحی شده
است. این سامانه یک رابط برنامه‌نویسی سازگار با \lr{OpenAI} ارائه می‌دهد که
امکان اتصال به چهار ارائه‌دهنده‌ی هوش مصنوعی شامل \lr{OpenAI}، \lr{Gemini}،
\lr{Groq} و \lr{OpenRouter} را فراهم می‌کند.

ویژگی‌های اصلی سامانه عبارت‌اند از: احراز هویت مبتنی بر
\term{رمز یک‌بارمصرف}{One-Time Password (OTP)}، مدیریت کلیدهای رابط
برنامه‌نویسی با محدودیت نرخ و لیست سفید آدرس \lr{IP}، سیستم کیف پول و صدور
صورت‌حساب، اندازه‌گیری مصرف \term{توکن}{Token} با استفاده از
\term{توکن‌ساز واقعی}{Real Tokenizer}، پشتیبانی از صورت‌حساب‌گذاری غیرتوکنی
(تصویر، صوت)، ارسال \term{گزارش}{Log} به سرویس‌های خارجی از طریق صف، و
گزارش‌دهی مصرف.

معماری سامانه بر \term{الگوی خط لوله}{Pipeline Pattern} استوار است که
پردازش هر درخواست را به چهارده مرحله‌ی مستقل تقسیم می‌کند. همچنین از
\term{استراتژی کش‌کنار}{Cache-Aside Strategy} برای کاهش بار پایگاه داده در
داده‌های پرمصرف استفاده شده است.

نتایج ارزیابی نشان می‌دهد که سامانه با موفقیت درخواست‌های هوش مصنوعی را
به ارائه‌دهندگان مختلف هدایت می‌کند، صورت‌حساب دقیق را بر اساس مصرف واقعی
محاسبه می‌کند و از طریق مجموعه‌ی آزمون‌های واحد و ویژگی پوشش مناسبی دارد.

\vspace{1cm}
\noindent\textbf{کلیدواژه‌ها:}
سامانه‌ی واسط، هوش مصنوعی، رابط برنامه‌نویسی یکپارچه، الگوی خط لوله،
مدیریت سرویس، صورت‌حساب، \lr{Laravel}، \lr{Redis}، \lr{OpenAI}

\cleardoublepage
